\subsection{从企业到国家：中美AI治理的宏观政治对照}

上一小节的产业对照揭示了一个更深层的问题：企业层面的开源与闭源选择，不是孤立的市场决策，而是嵌入在国家层面的AI治理战略之中。中美两国最高领导层对AI发展方向的不同判断，构成了企业AI获取模式的宏观政治条件。

\subsubsection{中国路径：新质生产力与开源普惠}

自2023年以来，中国最高领导层将人工智能明确纳入\textbf{新质生产力}的战略框架。在多个公开场合，习近平强调AI技术应服务于实体经济转型和高质量发展。2024年政府工作报告首次提出\textbf{AI+}行动计划，核心目标不是研发单一尖端模型，而是将AI能力扩散至制造业、农业、医疗、教育等各行业的毛细血管。

在这一战略方向下，开源AI生态获得了明确的政策支持：

\begin{itemize}
\item \textbf{产业逻辑}：如果AI是新质生产力的基础设施，那么它的获取门槛必须足够低。开源模型——如GLM、DeepSeek、通义千问的开源版本——使中小企业和县域经济无需高昂的API费用即可使用AI。这是\textbf{生产力逻辑}对\textbf{商业模式}的优先：国家战略层面关注的是AI的扩散速度和渗透率，而非单一企业的商业收入。
\item \textbf{数字主权逻辑}：本地部署的开源模型使中国企业在使用AI时不必依赖美国公司的API。在中美科技竞争背景下，这不仅是经济选择，也是供应链安全的战略选择。
\item \textbf{认知普惠逻辑}：与本文的认知民主化命题一致，中国政策文件中反复出现的AI赋能中小企业和基层治理的论述，实质上是在推动AI认知光谱的广泛爬升——让更多人从0-2级进入3-4级。
\end{itemize}

国务院发展研究中心和中国信息通信研究院的多份报告均指出，开源AI是实现AI普惠的关键路径。这与美国以安全管控和出口限制为主导的AI治理模式形成了系统性对照。

\subsubsection{美国路径：安全优先与技术遏制}

美国AI政策呈现出一种内在张力。一方面，硅谷的创新传统推崇开放和竞争；另一方面，自2023年以来，AI安全与对齐已成为华盛顿两党共识的核心议题。

拜登政府于2023年10月签署的AI行政令（Executive Order on Safe, Secure, and Trustworthy AI）要求最强大的AI模型在发布前进行安全测试并报告结果。这一行政令在客观上强化了大型闭源模型的优势——因为只有资源充足的大公司（如OpenAI、Google、Anthropic）才能承担合规成本。开源模型的分发者（如Meta的Llama系列）则面临更模糊的监管边界。

特朗普阵营则倾向于大规模放松AI监管，强调与中国在AI领域的全面竞争，但保留了技术出口管制的核心工具——特别是对华高端AI芯片的出口限制。这些限制虽然针对硬件而非软件，但对开源AI的生态也产生了间接影响：限制了中国开源模型在训练阶段的硬件资源，从而在一定程度上抑制了开源生态的竞争力。

布鲁金斯学会和兰德公司的多项研究报告指出，美国AI治理的核心焦虑在于\textbf{安全与竞争}的平衡：既担心AI被恶意行为者利用，又担心过度监管将创新让给中国。这种焦虑在实践中倾向于\textbf{加强管控}——无论是通过安全测试要求还是通过芯片出口限制。

\subsubsection{两种治理哲学的光谱映射}

将中美AI治理路径映射到本文的分析框架：

\begin{center}
\begin{tabular}{c|c|c}
\hline
维度 & 中国路径 & 美国路径 \\
\hline
核心叙事 & 新质生产力、AI普惠 & AI安全、技术竞争 \\
治理工具 & 产业政策+开源生态支持 & 安全行政令+出口管制 \\
企业AI获取模式 & 去中介化（本地部署）& 再中介化（API许可+PE审计） \\
对应本文光谱逻辑 & 降低光谱爬升门槛 & 强化守门人结构 \\
对开源的态度 & 政策鼓励+产业推动 & 安全担忧+模糊监管 \\
\hline
\end{tabular}
\end{center}

需要强调：这不是一个简单的二分——中国有AI安全法规（如生成式AI服务管理办法），美国也有活跃的开源AI社区（如Llama、Mistral的商业生态）。但两国顶层战略的\textbf{重心差异}是显著的：在扩散与管控之间，中国的政策天平更倾向扩散；在开放与安全之间，美国的政策天平更倾向安全。

\subsubsection{与本文理论框架的对接}

这一宏观政治对照强化了本文的核心论证：

\begin{enumerate}
\item \textbf{AI获取模式是政治选择，而非纯粹的市场结果}。开源vs闭源的选择不仅取决于企业层面的成本效益分析，更取决于国家战略是否将AI视为\textit{公共基础设施}（需要去中介化扩散）还是\textit{战略性商品}（需要管控和许可）。
\item \textbf{认知民主化的制度条件}：中国的开源AI政策在客观上降低了光谱爬升的门槛——不是出于抽象的民主理念，而是出于生产力扩散的实际需要。这提示了一个重要洞察：认知民主化的推动力不必然来自民主价值观，也可能来自\textit{发展主义的实用逻辑}。
\item \textbf{开源作为地缘技术选择}：在中美科技竞争背景下，开源不仅是认知民主化的工具，也是技术主权的保障。一个依赖闭源API的AI生态系统，在地缘政治紧张时面临更大的断供风险。
\end{enumerate}

这一分析也为本文的政策含义增加了国际维度：认知民主化不仅是一国之内的平等问题，也是国际体系中的权力问题。当AI成为新时代的基础设施时，谁控制基础设施的入口，谁就控制了基础设施使用者的认知进化路径。
